% Compile with XeLaTeX (twice).
\documentclass[11pt,UTF8,fontset=windows]{ctexart}
\usepackage[letterpaper,margin=0.85in]{geometry}
\usepackage{amsmath,amssymb,fancyhdr,xcolor,listings,pdfpages}
\usepackage[hidelinks]{hyperref}
\pagestyle{fancy}\fancyhf{}
\fancyhead[L]{EECS 453 / Homework 1}\fancyhead[R]{解答与推导}
\fancyfoot[C]{\thepage}\setlength{\headheight}{15pt}
\setlength{\parindent}{0pt}\setlength{\parskip}{4pt}\linespread{1.12}
\ctexset{section={beforeskip=10pt,afterskip=7pt},subsection={beforeskip=8pt,afterskip=4pt}}
\lstset{basicstyle=\ttfamily\small,breaklines=true,columns=fullflexible,frame=single,keepspaces=true,showstringspaces=false}
\newcommand{\R}{\mathbb R}
\begin{document}
\begin{center}{\LARGE Homework 1 解答}\\[5pt]线性代数、概率、最小二乘与 NumPy\end{center}
\section*{Problem 1：线性代数复习}
\subsection*{(a) 证明转置恒等式}
设 $A$ 的第 $i$ 行为 $a_i^T$，其中 $a_i\in\R^n$ 是列向量。于是
\[
 A=\begin{bmatrix}a_1^T\\\vdots\\a_n^T\end{bmatrix},\qquad
 Av=\begin{bmatrix}a_1^Tv\\\vdots\\a_n^Tv\end{bmatrix},\qquad
 A^T=\begin{bmatrix}a_1&\cdots&a_n\end{bmatrix}.
\]
因为 $a_i^Tv$ 为实数，$a_i^Tv=v^Ta_i$，所以
\[
 (Av)^T=\begin{bmatrix}a_1^Tv&\cdots&a_n^Tv\end{bmatrix}
 =\begin{bmatrix}v^Ta_1&\cdots&v^Ta_n\end{bmatrix}
 =\boxed{v^TA^T}.
\]
\subsection*{(b) 链式法则}
这里求的是复合函数 $f(x(w),y(w))$ 对 $w$ 的导数。首先
\[
 f_x=2x+2y,\quad f_y=2x-2y,\quad x'=2w-3,\quad y'=4w^3-6w+1.
\]
由链式法则，
\begin{align*}
 \frac{df}{dw}
 &=(2x+2y)(2w-3)+(2x-2y)(4w^3-6w+1)\\
 &=2(w^4-2w^2-2w)(2w-3)
   +2(-w^4+4w^2-4w)(4w^3-6w+1)\\
 &=\boxed{-8w^7+48w^5-40w^4-56w^3+60w^2+4w}.
\end{align*}
\subsection*{(c) 正交变换保持平方范数}
正交矩阵满足 $U^TU=I$，因此
\[
 \|Ux\|_2^2=(Ux)^T(Ux)=x^TU^TUx=x^Tx=\boxed{\|x\|_2^2}.
\]
\subsection*{(d) 二次型的梯度}
展开为 $f(x)=\sum_{i=1}^n\sum_{j=1}^n A_{ij}x_ix_j$。对 $x_k$ 求偏导：
\[
 \frac{\partial f}{\partial x_k}
 =\sum_{j=1}^n A_{kj}x_j+\sum_{i=1}^n A_{ik}x_i.
\]
第一项来自 $x_i$ 的导数，第二项来自 $x_j$ 的导数；当 $i=j=k$ 时，两项合起来正好是 $2A_{kk}x_k$。
故 $\nabla_x f=(A+A^T)x$。利用 $A=A^T$，得到
\[\boxed{\nabla_x f(x)=2Ax}.\]
\newpage
\section*{Problem 2：概率复习}
\subsection*{(a) 恰好匹配 $k$ 个号码}
固定购买的彩票，令 $K$ 表示它与中奖号码的匹配数。中奖号码是从 51 个号码中等概率选取的 6 元子集，共有 $\binom{51}{6}$ 种。

若恰好匹配 $k$ 个，需要从彩票的 6 个号码中选出 $k$ 个，再从其余 45 个号码中选出 $6-k$ 个。两个选择共同且唯一地确定中奖集合，所以
\[
 \boxed{\Pr(K=k)=\frac{\binom{6}{k}\binom{45}{6-k}}{\binom{51}{6}},
 \qquad k=0,1,\ldots,6.}
\]
这是超几何分布。注意号码不计顺序，因此使用组合数而非排列数。
\subsection*{(b) 机器说谎的概率}
按通常约定，骰子为相互独立的公平六面骰，掷骰结果也独立于抽奖。令 $L$ 表示机器说谎。

当 $K=0$ 时，机器必定说实话，即 $\Pr(L\mid K=0)=0$。
当 $K=k\geq1$ 时，$k$ 个骰子的点数之和至少为 $k$。和等于 $k$ 当且仅当每个骰子都是 1，因此
\[
 \Pr(L\mid K=k)=\left(\frac16\right)^k,\qquad k=1,\ldots,6.
\]
应用全概率公式，
\begin{align*}
 \Pr(L)&=\sum_{k=0}^6\Pr(L\mid K=k)\Pr(K=k)\\
 &=\boxed{\frac{1}{\binom{51}{6}}
 \sum_{k=1}^6\binom{6}{k}\binom{45}{6-k}\frac{1}{6^k}}\\
 &=\frac{11992137545}{168049873152}
 \approx 0.071360586712.
\end{align*}
所以机器说谎的概率约为 $\boxed{7.1361\%}$。求和从 $k=1$ 开始，不能把 $k=0$ 的情况也计作说谎。
\newpage
\section*{Problem 3：线性最小二乘}
\subsection*{(a) 在特征向量坐标系中表示目标函数}
由实对称矩阵的谱分解，$A=Q\Lambda Q^T$，其中 $Q^TQ=QQ^T=I$。定义
\[\boxed{\widetilde x=Q^Tx,\qquad \widetilde y=Q^Ty}.\]
因为 $y=Q\widetilde y$，故
\[
 Ax-y=Q\Lambda Q^Tx-QQ^Ty=Q(\Lambda\widetilde x-\widetilde y).
\]
根据 Problem 1(c) 中的正交不变性，
\[\boxed{\|Ax-y\|_2^2=\|\Lambda\widetilde x-\widetilde y\|_2^2}.\]
且 $x\leftrightarrow\widetilde x$ 是一一对应的坐标变换，逆变换为 $x=Q\widetilde x$，因此在两组坐标下最小化是等价的。
\subsection*{(b) 伪逆给出一个最小化解}
令 $I_+=\{i:\lambda_i\ne0\}$，$I_0=\{i:\lambda_i=0\}$。这里 $I_+$ 仅表示非零指标，并不要求特征值为正。逐坐标展开：
\[
 \|\Lambda\widetilde x-\widetilde y\|_2^2
 =\sum_{i\in I_+}(\lambda_i\widetilde x_i-\widetilde y_i)^2
  +\sum_{i\in I_0}\widetilde y_i^2.
\]
第二项与 $\widetilde x$ 无关；第一项非负，且取
$\widetilde x_i=\widetilde y_i/\lambda_i$ 时每个平方项都为零。
对 $i\in I_0$，$\widetilde x_i$ 可以任意选择。特别地，取它们为零，则
\[
 \widetilde x^*=\Lambda^\dagger\widetilde y,
 \qquad (\Lambda^\dagger)_{ii}=\begin{cases}1/\lambda_i,&\lambda_i\ne0,\\0,&\lambda_i=0,\end{cases}
\]
非对角元素全为零。变回原坐标得到
\[\boxed{x^*=Q\Lambda^\dagger Q^Ty},\qquad
 \boxed{\min_x f(x)=\sum_{i\in I_0}\widetilde y_i^2}.\]
\textbf{解的非唯一性：}由于 $A$ 不可逆，至少一个零特征值存在。所有最小化解为
\[
 x=x^*+\sum_{i\in I_0}c_iq_i,\qquad c_i\in\R.
\]
这些额外项位于 $\ker A$，不改变 $Ax$。又由于它们与 $x^*$ 正交，
$\|x\|_2^2=\|x^*\|_2^2+\sum_{i\in I_0}c_i^2$，所以 $x^*$ 是唯一的最小范数最小二乘解。
\newpage
\subsection*{(c) 加分题：一般矩阵的 SVD 解法}
设 $A\in\R^{m\times n}$，$y\in\R^m$，$\operatorname{rank}(A)=r$。取完整奇异值分解
\[
 A=U\Sigma V^T,\quad U\in\R^{m\times m},\quad
 V\in\R^{n\times n},\quad\Sigma\in\R^{m\times n},
\]
其中 $U,V$ 正交，$\Sigma$ 的前 $r$ 个对角元素为 $\sigma_i>0$，其余元素为零。定义
$z=V^Tx\in\R^n$，$b=U^Ty\in\R^m$。于是
\begin{align*}
 \|Ax-y\|_2^2
 &=\|U(\Sigma z-b)\|_2^2=\|\Sigma z-b\|_2^2\\
 &=\sum_{i=1}^r(\sigma_i z_i-b_i)^2+\sum_{i=r+1}^m b_i^2.
\end{align*}
这同时适用于高矩阵、宽矩阵和秩亏矩阵。后一个求和为无法通过选择 $x$ 消去的常数；令
\[
 z_i^*=\begin{cases}b_i/\sigma_i,&1\le i\le r,\\0,&r<i\le n,\end{cases}
\]
便使第一个求和达到零。定义 $\Sigma^\dagger\in\R^{n\times m}$，它的前 $r$ 个对角元素为 $1/\sigma_i$，其余为零。则
\[
 \boxed{x^*=V\Sigma^\dagger U^Ty=A^\dagger y},\qquad
 \boxed{\min_x\|Ax-y\|_2^2=\sum_{i=r+1}^m b_i^2}.
\]
若 $v_i$ 是 $V$ 的第 $i$ 列，所有最小化解为
\[
 x=x^*+\sum_{i=r+1}^n c_iv_i,\qquad c_i\in\R.
\]
由 $V$ 的正交性，$x^*$ 仍是唯一的最小范数最小二乘解。当 $r=n$ 时最小化解唯一；当 $r<n$ 时不唯一。
当 $r=0$ 时，$A=0$，所有 $x$ 都是最小化解，而 $x^*=0$ 是最小范数解。空求和按零处理，因此上述公式也覆盖此情形。
\newpage
\section*{Problem 4：NumPy 向量化}
以下答案均不使用循环或列表推导式。附录是完成后的 notebook 的代码单元与实际运行输出；原 notebook 的参考检查函数保持不变。
\subsection*{(a) 各行的 $\ell_1$ 范数}
\begin{lstlisting}[language=Python]
def row_l1_norms(A):
    return np.abs(A).sum(axis=1)
\end{lstlisting}
对 $A\in\R^{N\times M}$，第 $i$ 个输出为 $x_i=\sum_{j=1}^M|A_{ij}|$。
绝对值逐元素计算，随后沿列方向（\texttt{axis=1}）求和，输出形状为 \texttt{(N,)}。函数体只有一行。
\subsection*{(b) 为每个住址分配最近基站}
\begin{lstlisting}[language=Python]
def closest_tower(A, B):
    a2 = (A**2).sum(axis=1)[None, :]
    b2 = (B**2).sum(axis=1)[:, None]
    sq_dists = b2 + a2 - 2 * B @ A.T
    return np.argmin(sq_dists, axis=1)
\end{lstlisting}
函数体恰好四行。假设有至少一个基站（$M\ge1$）。各个中间量的形状如下：
\begin{center}
\begin{tabular}{lll}
\hline
变量 & 形状 & 含义\\\hline
\texttt{a2} & \texttt{(1, M)} & 各基站坐标的平方范数\\
\texttt{b2} & \texttt{(N, 1)} & 各住址坐标的平方范数\\
\texttt{B @ A.T} & \texttt{(N, M)} & 所有住址与基站的内积\\
\texttt{sq\_dists} & \texttt{(N, M)} & 所有住址与基站的平方距离\\
输出 & \texttt{(N,)} & 各住址的最近基站索引\\\hline
\end{tabular}
\end{center}
广播相加并减去内积项后，第 $(i,j)$ 个元素为
\[
 \texttt{sq\_dists}_{ij}=\|B_{i,:}\|_2^2+\|A_{j,:}\|_2^2
 -2B_{i,:}A_{j,:}^T=\|B_{i,:}-A_{j,:}\|_2^2.
\]
平方根严格递增，因此无需开平方就能找最近基站。沿 \texttt{axis=1} 在每行的 $M$ 个基站中取最小值索引。距离相同时，\texttt{np.argmin} 返回第一个最小值的索引。

原代码有两处错误：内积前应为负号；\texttt{argmin} 应沿 \texttt{axis=1}，才能为每个住址返回一个索引。
\clearpage
\section*{附录：完成后的 notebook}
以下按原顺序列出全部代码单元及运行输出。原题的检查单元包含参考循环，仅用于验证，不属于答案函数。
\subsection*{In [1]}
\begin{lstlisting}[language=Python]
import numpy as np

rng = np.random.default_rng(453)
\end{lstlisting}
\subsection*{In [2]}
\begin{lstlisting}[language=Python]
def row_l1_norms(A):
    return np.abs(A).sum(axis=1)
\end{lstlisting}
\subsection*{In [3]}
\begin{lstlisting}[language=Python]
# --- check (a) ---
A_test = rng.normal(size=(5, 4))

def _ref_row_l1(A):            # slow reference, for checking only -- do not use in your answer
    return np.array([np.abs(A[i]).sum() for i in range(A.shape[0])])

out = row_l1_norms(A_test)
assert out.shape == (5,), f"expected shape (5,), got {out.shape}"
assert np.allclose(out, _ref_row_l1(A_test)), "values do not match"
print("(a) looks correct")
\end{lstlisting}
\textbf{实际输出}
\begin{lstlisting}
(a) looks correct
\end{lstlisting}
\newpage
\subsection*{In [4]}
\begin{lstlisting}[language=Python]
def closest_tower(A, B):
    a2 = (A**2).sum(axis=1)[None, :]
    b2 = (B**2).sum(axis=1)[:, None]
    sq_dists = b2 + a2 - 2 * B @ A.T
    return np.argmin(sq_dists, axis=1)
\end{lstlisting}
\subsection*{In [5]}
\begin{lstlisting}[language=Python]
# --- check (b) ---
A_towers = rng.normal(size=(7, 3))     # M = 7 towers,  D = 3
B_homes  = rng.normal(size=(11, 3))    # N = 11 homes

def _ref_closest(A, B):        # slow reference, for checking only -- do not use in your answer
    return np.array([int(np.argmin(((B[i] - A) ** 2).sum(axis=1))) for i in range(B.shape[0])])

out = closest_tower(A_towers, B_homes)
assert out.shape == (B_homes.shape[0],), f"expected shape {(B_homes.shape[0],)}, got {out.shape}"
assert np.array_equal(out, _ref_closest(A_towers, B_homes)), "indices are not the nearest towers"
print("(b) looks correct")
\end{lstlisting}
\textbf{实际输出}
\begin{lstlisting}
(b) looks correct
\end{lstlisting}
\end{document}
